\begin{frame}{讲解：统一状态视角到底统一了什么}
  \Core{这一讲不是把所有缓存混成一个词，而是把不同对象放到同一分析坐标里。}
  \begin{itemize}
    \item KV、检索知识块、会话记忆都具有跨时间存在、可能复用、需要回收这三个共同特征。
    \item 它们的差别在于生成方式、正确性要求、搬运代价和复用条件不同。
    \item 所谓“统一状态管理”，不是让它们实现完全一样，而是用同一套语言讨论驻留、迁移、命名和价值。
  \end{itemize}
\end{frame}

\begin{frame}{讲解：第六讲应形成的研究意识}
  \Core{状态一旦进入主路径，系统问题就不再只是计算快慢，而变成“对象如何持续存在”。}
  \begin{itemize}
    \item 只谈计算资源而不谈状态落点，很多机制会在真实系统里失效。
    \item 只谈潜在复用而不谈搬运与一致性，很多收益会停留在理想叙述里。
    \item 因此这一讲要让学生意识到：推理系统正在从执行系统走向状态化在线系统。
  \end{itemize}
\end{frame}

\begin{frame}{例子：为什么 DistServe 要把 prefill 和 decode 拆开（1）}
  \Core{状态管理不是“把对象存起来”，而是决定对象该在什么地方、什么时机被持有和搬运。}
\begin{block}{单机一体化路径}
    prefill 和 decode 共用一套资源，简单但容易让两类阶段互相干扰。
  \end{block}
\end{frame}

\begin{frame}{例子：为什么 DistServe 要把 prefill 和 decode 拆开（2）}
  \Core{状态管理不是“把对象存起来”，而是决定对象该在什么地方、什么时机被持有和搬运。}
\begin{block}{阶段解耦路径}
    prefill 更靠近高吞吐算力，decode 更靠近低抖动稳态资源，状态搬运因此成为主问题。
  \end{block}
\end{frame}

\begin{frame}{例子：为什么 DistServe 要把 prefill 和 decode 拆开（3）}
  \Core{状态管理不是“把对象存起来”，而是决定对象该在什么地方、什么时机被持有和搬运。}
\begin{block}{课堂结论}
    一旦阶段被拆开，系统研究的重点就从“存没存”转到“搬得动、对不对、值不值”。
  \end{block}
\end{frame}

\begin{frame}{例子：RAGCache 为什么让状态视角进一步扩张（1）}
  \Core{课程中的“状态”不只包括 KV，还包括检索知识块、会话记忆和调度元数据。}
\begin{block}{传统理解}
    缓存常被理解成底层实现优化，像一个被动的加速器。
  \end{block}
\end{frame}

\begin{frame}{例子：RAGCache 为什么让状态视角进一步扩张（2）}
  \Core{课程中的“状态”不只包括 KV，还包括检索知识块、会话记忆和调度元数据。}
\begin{block}{统一状态视角}
    检索结果、前缀块和会话历史都可能跨请求复用，因此它们都值得被当成一等状态对象建模。
  \end{block}
\end{frame}

\begin{frame}{例子：RAGCache 为什么让状态视角进一步扩张（3）}
  \Core{课程中的“状态”不只包括 KV，还包括检索知识块、会话记忆和调度元数据。}
\begin{block}{系统含义}
    一旦对象具有跨时间和跨请求价值，系统就必须管理它的命名、引用、迁移和回收边界。
  \end{block}
\end{frame}

\begin{frame}{代码例子：Sequence 的序列化为什么值得课堂关注（1）}
  \Core{状态管理课程里的代码例子，应让学生看到对象不只是运行时临时变量，还需要可传递与可恢复。}
\begin{block}{代码片段}
    {\ttfamily\small
    \detokenize{def __getstate__(self): ...}\par
    \detokenize{def __setstate__(self, state): ...}\par
    }
  \end{block}
\end{frame}

\begin{frame}{代码例子：Sequence 的序列化为什么值得课堂关注（2）}
  \Core{状态管理课程里的代码例子，应让学生看到对象不只是运行时临时变量，还需要可传递与可恢复。}
\begin{block}{课堂应追问}
    哪些字段如果丢失，请求将无法继续？哪些字段如果错传，会直接造成状态不一致？
  \end{block}
\end{frame}

\begin{frame}{代码例子：Sequence 的序列化为什么值得课堂关注（3）}
  \Core{状态管理课程里的代码例子，应让学生看到对象不只是运行时临时变量，还需要可传递与可恢复。}
\begin{block}{系统含义}
    只要对象可能跨模块、跨节点或跨阶段流动，序列化语义就是系统正确性的一部分。
  \end{block}
\end{frame}

\begin{frame}{代码例子：统一状态管理器会暴露哪些接口（1）}
  \Core{课程项目里一个常见扩展，是想象如果把 KV、知识块和会话记忆统一管理，需要哪些最小接口。}
\begin{block}{接口草图}
    {\ttfamily\small
    \detokenize{state_id = state_store.put(obj, metadata)}\par
    \detokenize{state_store.pin(state_id, device="gpu")}\par
    \detokenize{state_store.evict(policy="lru") }\par
    }
  \end{block}
\end{frame}

\begin{frame}{代码例子：统一状态管理器会暴露哪些接口（2）}
  \Core{课程项目里一个常见扩展，是想象如果把 KV、知识块和会话记忆统一管理，需要哪些最小接口。}
\begin{block}{课堂应追问}
    这些接口背后分别对应命名、驻留、迁移还是回收问题？
  \end{block}
\end{frame}

\begin{frame}{代码例子：统一状态管理器会暴露哪些接口（3）}
  \Core{课程项目里一个常见扩展，是想象如果把 KV、知识块和会话记忆统一管理，需要哪些最小接口。}
\begin{block}{结论}
    统一状态管理的难点不是接口名字，而是不同状态对象的价值函数和一致性需求并不相同。
  \end{block}
\end{frame}

\begin{frame}{例子：状态搬运为什么经常吞掉理论收益（1）}
  \Core{很多机制论文的理论收益很强，但一旦进入真实系统，搬运与同步代价就会迅速显形。}
\begin{block}{理想世界}
    把更适合 prefill 的资源和更适合 decode 的资源拆开，看起来可以两头占优。
  \end{block}
\end{frame}

\begin{frame}{例子：状态搬运为什么经常吞掉理论收益（2）}
  \Core{很多机制论文的理论收益很强，但一旦进入真实系统，搬运与同步代价就会迅速显形。}
\begin{block}{真实世界}
    KV 或知识对象一旦要跨节点移动，带宽、延迟、重组和一致性全都会冒出来。
  \end{block}
\end{frame}

\begin{frame}{例子：状态搬运为什么经常吞掉理论收益（3）}
  \Core{很多机制论文的理论收益很强，但一旦进入真实系统，搬运与同步代价就会迅速显形。}
\begin{block}{教学重点}
    讲状态复用时必须同时讲状态搬运，否则学生会高估机制落地的直接性。
  \end{block}
\end{frame}

\begin{frame}{课堂练习：把三类状态放到同一张表里（1）}
  \Core{学生需要学会比较不同状态对象，而不是只会孤立讨论 KV。}
\begin{block}{练习一}
    比较 KV、RAG 检索块、会话记忆三者在生成方式、驻留时间和复用条件上的差异。
  \end{block}
\end{frame}

\begin{frame}{课堂练习：把三类状态放到同一张表里（2）}
  \Core{学生需要学会比较不同状态对象，而不是只会孤立讨论 KV。}
\begin{block}{练习二}
    说明哪一类状态最适合先做 observe-only 探针，哪一类最需要严格正确性边界。
  \end{block}
\end{frame}

\begin{frame}{课堂练习：把三类状态放到同一张表里（3）}
  \Core{学生需要学会比较不同状态对象，而不是只会孤立讨论 KV。}
\begin{block}{练习三}
    用一句话回答：为什么第 6 讲的关键词不是“缓存”，而是“状态组织”？
  \end{block}
\end{frame}
